% !TEX root = ../../../main.tex

\subsection{句子视频识别任务定位}

本项目中的句子视频识别并不等同于开放域连续手语翻译，而是在有限词表和受限场景下对短视频识别链路进行验证。
系统希望完成的是一条可接入 HearBridge 的工程流程：用户选择本地手语视频，识别服务输出词级候选结果，服务端组织中文映射，并进一步通过语义修正得到更自然的展示结果。

与 6.1 节中的实时单词识别相比，句子视频识别的输入形式、处理目标和输出结构都发生了变化。
实时单词识别面对的是移动端摄像头中的连续图像流，重点在于对当前训练动作进行即时反馈；
句子视频识别面对的是用户已经录制或选择的一段短视频，重点在于从完整视频中定位若干可能连续出现的词级动作，并按时间顺序组织为 Gloss 序列。
因此，本节所讨论的“句子”并不是自然语言意义上的完整句法生成结果，而是由有限词表中的多个词级识别结果构成的短语或电报式表达。
后续的中文直译和语义自然化均建立在该 Gloss 序列之上，不能替代视觉识别模块本身。

从系统职责看，句子视频识别链路被划分为视觉识别层、业务编排层和结果展示层。
Python 识别服务负责视频解码、MediaPipe 关键点检测、词级模型推理和候选词段输出；
Spring Boot 后端负责接收移动端上传请求、调用 Python 服务、保存识别记录，并在需要时调用 DeepSeek 对展示文本进行语义自然化；
HarmonyOS 移动端负责选择或录制视频，并向用户展示原始 Gloss、中文直译和自然化文本。
这种分层方式可以保证视觉识别结果被完整保留，同时允许展示层在不覆盖原始结果的前提下提高可读性。

表~\ref{tab:chapter06-sentence-task-scope} 给出了句子视频识别任务的边界与输出层次。
其中，原始识别结果主要用于评价模型与后处理算法，中文直译用于帮助普通用户理解 Gloss 含义，自然化文本则作为辅助展示结果。
后续 6.5 节中的帧级检测缓存、滑动窗口预测、词段合并和 DeepSeek 候选修正流程，分别如图~\ref{fig:chapter06-frame-cache-flow}、图~\ref{fig:chapter06-sliding-window}、图~\ref{fig:chapter06-window-predict-to-gloss} 和图~\ref{fig:chapter06-deepseek-candidate-correction} 所示。
这些流程共同构成了从上传视频到最终展示文本的完整识别链路。

\begin{longtable}{L{0.20\textwidth} L{0.28\textwidth} L{0.44\textwidth}}
  \caption{句子视频识别任务边界说明表}
  \label{tab:chapter06-sentence-task-scope}\\
  \toprule
  层次 & 输入输出 & 任务边界 \\
  \midrule
  视频输入 & 用户上传的短视频文件 & 处理离线视频，不要求摄像头实时连续返回结果 \\
  视觉识别 & 视频帧、关键点与词级模型 & 输出词段、置信度和 Top-K 候选，不直接生成最终自然语言句子 \\
  序列组织 & 按时间排列的词段结果 & 通过过滤、合并和重叠抑制形成原始 Gloss 序列 \\
  中文映射 & Gloss 标签与中文词义表 & 将 rawSequence 映射为 rawDisplayZh 和 rawTextZh \\
  展示修正 & 原始序列与候选集合 & 只对展示文本进行自然化，保留原始视觉识别结果 \\
  \bottomrule
\end{longtable}

由表~\ref{tab:chapter06-sentence-task-scope} 可知，本项目将句子视频识别定位为“有限词表词段识别 + 展示层语义组织”的工程验证任务。
这种定位一方面避免将小词表实验夸大为开放域手语翻译系统，另一方面也便于将识别结果接入 HearBridge 的移动端、后端和管理端业务流程。
在后续实验分析中，模型性能主要通过词级分类效果、拼接句子视频的序列匹配情况以及语义修正前后的展示效果进行评价。
